% 数学ノート集・共通テンプレート v2（積分論2026年度第1部の版面を参考）
% A5判。章内で定義・定理・例・注意等の番号を共通化。
% コンパイル： platex -interaction=nonstopmode math_note_template_v2.tex (2回)
%              dvipdfmx math_note_template_v2.dvi
\documentclass[a5paper,10pt,onecolumn]{jsarticle}
\usepackage[dvipdfmx,top=22mm,bottom=20mm,inner=19mm,outer=18mm,headheight=13pt,headsep=6mm]{geometry}
\usepackage{amsmath,amssymb,amsthm,mathtools}
\usepackage{enumitem,fancyhdr,titlesec}
\usepackage[dvipdfmx,bookmarks=true,colorlinks=false]{hyperref}
\usepackage{pxjahyper}

% ----- 差し替える文書情報 -----
\newcommand{\NoteTitle}{曲率テンソルとJacobi場}
\newcommand{\NoteSeries}{数学ノート集}

% ----- 文字・組版 -----
% 参考PDFの質素な明朝＋Computer Modernの印象を引き継ぎながら，
% A5判で窮屈にならないよう本文の行間を少し確保する．
\linespread{1.13}
\setlength{\parindent}{1zw}
\setlength{\parskip}{0pt}
\setlength{\abovedisplayskip}{6pt plus 2pt minus 2pt}
\setlength{\belowdisplayskip}{6pt plus 2pt minus 2pt}
\setlength{\abovedisplayshortskip}{3pt plus 2pt minus 1pt}
\setlength{\belowdisplayshortskip}{4pt plus 2pt minus 1pt}
\setlist[enumerate]{topsep=3pt,itemsep=1pt,parsep=0pt,leftmargin=2.4zw}
\titleformat{\section}{\normalfont\normalsize\bfseries}{\thesection}{0.8zw}{}
\titlespacing*{\section}{0pt}{2.0ex plus .4ex}{.9ex}
\titleformat{\subsection}{\normalfont\normalsize\bfseries}{\thesubsection}{0.7zw}{}
\titlespacing*{\subsection}{0pt}{1.5ex plus .3ex}{.6ex}
\raggedbottom

% ----- 柱とページ番号：ページ上部、罫線なし -----
\fancypagestyle{mathnote}{%
  \fancyhf{}%
  \fancyhead[L]{\footnotesize \NoteSeries\ （\NoteTitle）}%
  \fancyhead[R]{\footnotesize \thepage}%
  \renewcommand{\headrulewidth}{0pt}%
  \renewcommand{\footrulewidth}{0pt}%
}
\pagestyle{mathnote}

% ----- 節内通し番号の定義・定理群 -----
\newtheoremstyle{mathnotetheorem}
 {4.5pt}{4.5pt}{\normalfont}{}%
 {\bfseries}{}{0.65zw}%
 {\thmname{#1}\thmnumber{ #2}\thmnote{（#3）}}
\theoremstyle{mathnotetheorem}
\newtheorem{theorem}{定理}[section]
\newtheorem{definition}[theorem]{定義}
\newtheorem{lemma}[theorem]{補題}
\newtheorem{proposition}[theorem]{命題}
\newtheorem{corollary}[theorem]{系}
\newtheorem{example}[theorem]{例}
\newtheorem{remark}[theorem]{注意}
\renewcommand{\proofname}{証明}
\renewcommand{\qedsymbol}{$\square$}
\makeatletter
\renewenvironment{proof}[1][\proofname]{%
  \par\pushQED{\qed}\normalfont\topsep4\p@\@plus2\p@\relax
  \trivlist\item[\hskip\labelsep\bfseries #1]\ignorespaces
}{\popQED\endtrivlist\@endpefalse}
\makeatother

% ----- 使うノートで適宜拡張する記号 -----
\newcommand{\R}{\mathbb{R}}
\newcommand{\norm}[1]{\lVert #1\rVert}
\newcommand{\abs}[1]{\lvert #1\rvert}

\begin{document}
\section*{読了後にできること}
本稿を読むための前提は，Riemann曲率テンソル，測地線とJacobi方程式．
本稿の具体的な到達課題は，定曲率空間でJacobi場を解き，曲率と測地線の収束・発散の関係を説明する．


\section{Riemann曲率テンソル}
本稿では $(M,g)$ をRiemann多様体，$\nabla$ を
Levi-Civita接続とする．
\begin{definition}
\[
 R(X,Y)Z=\nabla_X\nabla_YZ-\nabla_Y\nabla_XZ
 -\nabla_{[X,Y]}Z
\]
を曲率テンソルといい，
$R(X,Y,Z,W)=g(R(X,Y)Z,W)$ と置く．
\end{definition}
\begin{proposition}[テンソル性]
$R(X,Y)Z$ は $X,Y,Z$ について $C^\infty(M)$-線形であり，
$R(X,Y)=-R(Y,X)$ である．
\end{proposition}
\begin{proof}
$X$ を $fX$ に変えたときの余分な $Y(f)\nabla_XZ$ は
$[fX,Y]=f[X,Y]-Y(f)X$ からの項と相殺する．
$Y$ については反対称性を用いればよい．
$Z$ を $fZ$ に変えると生じる二階微分項は
$X(Yf)-Y(Xf)-[X,Y]f=0$ により相殺し，
一階の項も二つずつ相殺する．
反対称性は定義を入れ替えれば明らかである．
\end{proof}
\section{曲率の対称性}
\begin{theorem}[第一Bianchi恒等式]\label{thm:bianchi}
\[
 R(X,Y)Z+R(Y,Z)X+R(Z,X)Y=0.
\]
\end{theorem}
\begin{proof}
全てがテンソル的なので一点の近傍の局所座標で
$X,Y,Z$ を座標ベクトル場として証明すればよい．
それらの括弧は零で，捩率零から
$\nabla_XY=\nabla_YX$ である．
巡回和を展開すると
\[
 \nabla_X\nabla_YZ-\nabla_X\nabla_ZY
 +\nabla_Y\nabla_ZX-\nabla_Y\nabla_XZ
 +\nabla_Z\nabla_XY-\nabla_Z\nabla_YX=0
\]
となる．
\end{proof}
\begin{proposition}[計量対称性]\label{prop:curvsym}
\[
 \begin{aligned}
 R(X,Y,Z,W)&=-R(Y,X,Z,W),\\
 R(X,Y,Z,W)&=-R(X,Y,W,Z),\\
 R(X,Y,Z,W)&=R(Z,W,X,Y).
 \end{aligned}
\]
\end{proposition}
\begin{proof}
第1式は定義による．
計量保存を二度適用し
$[X,Y]g(Z,W)=X(Yg(Z,W))-Y(Xg(Z,W))$
を展開すると
$g(R(X,Y)Z,W)+g(Z,R(X,Y)W)=0$ を得るので
第2式が成立する．
第1，2式およびBianchi恒等式を用い，
\[
 R(X,Y,Z,W)+R(Y,Z,X,W)+R(Z,X,Y,W)=0
\]
と，$W$ を先頭へ移した同種の式を組み合わせると
$R(X,Y,Z,W)-R(Z,W,X,Y)=0$ を得る．
この最後の等式は，二つのBianchi恒等式の
項を交代性で並べ替え，
共通する項を消去したものである．
\end{proof}
\section{断面曲率}
\begin{definition}
$u,v\in T_pM$ が線形独立であるとき，
それらが張る平面 $\sigma$ の断面曲率を
\[
 K(\sigma)=
 \frac{g(R(u,v)v,u)}
 {g(u,u)g(v,v)-g(u,v)^2}
\]
で定義する．
\end{definition}
\begin{proposition}
断面曲率は平面 $\sigma$ の基底 $u,v$ の選択に依存しない．
\end{proposition}
\begin{proof}
$u,v$ から $u',v'$ への基底変換行列を
$A\in\operatorname{GL}_2(\mathbb R)$ とする．
$R$ の第1・第2引数の交代性と第3・第4引数の交代性から
分子は $(\det A)^2$ 倍となる．
分母はGram行列の行列式であり，
$G\mapsto A^TGA$ により同じく $(\det A)^2$ 倍となる．
比は不変である．
\end{proof}
\section{測地線の変分とJacobi場}
\begin{definition}
$\gamma$ を測地線とする．滑らかな写像
$F:(-\varepsilon,\varepsilon)\times I\to M$ で
各 $s$ 固定の曲線 $\gamma_s(t)=F(s,t)$ が測地線，
$\gamma_0=\gamma$ となるものを測地線変分という．
\[
 J(t)=\left.\frac{\partial F}{\partial s}\right|_{s=0}
\]
をその変分ベクトル場という．
\end{definition}
\begin{theorem}[Jacobi方程式]\label{thm:jacobi}
測地線変分の変分ベクトル場 $J$ は
\[
 \nabla_{\dot\gamma}\nabla_{\dot\gamma}J
 +R(J,\dot\gamma)\dot\gamma=0
\]
を満たす．
\end{theorem}
\begin{proof}
$T=\partial_tF$，$S=\partial_sF$ と置く．
パラメータ方向は可換なので捩率零から
$\nabla_TS=\nabla_ST$ である．
また各 $\gamma_s$ は測地線なので
$\nabla_TT=0$．
曲率の定義から
\[
 \nabla_T\nabla_TS
 =\nabla_T\nabla_ST
 =\nabla_S\nabla_TT+R(T,S)T
 =-R(S,T)T.
\]
$s=0$ へ制限すると主張の式を得る．
\end{proof}
\begin{theorem}[Jacobi場の初期値問題]
任意の $a,b\in T_{\gamma(0)}M$ に対し，
$J(0)=a$，
$\nabla_{\dot\gamma}J(0)=b$ を満たす
Jacobi場が一意に存在する．
\end{theorem}
\begin{proof}
$\gamma$ に沿う平行正規直交フレームを取ると，
Jacobi方程式は滑らかな係数行列を持つ
線形二階連立常微分方程式
$j''(t)+A(t)j(t)=0$ になる．
一階の線形系へ直せるので，
線形ODEの存在一意性により任意の初期値に対して
全ての $\gamma$ の定義域で唯一解を持つ．
\end{proof}
\begin{proposition}[指数写像の微分とJacobi場]
$p\in M$，$v,w\in T_pM$ とする．
$t\mapsto\exp_p(tv)$ に沿って
$J(0)=0$，
$\nabla_{\dot\gamma}J(0)=w$ となるJacobi場 $J$ を取る．
このとき
\[
 J(t)=(d\exp_p)_{tv}(tw)
\]
が成立する（定義域に含まれる $t$ について）．
\end{proposition}
\begin{proof}
測地線変分
$F(s,t)=\exp_p(t(v+sw))$ を取ると，
その変分場は定義から右辺となる．
$F(s,0)=p$ なので $J(0)=0$．
$\partial_tF(s,0)=v+sw$ を $s$ で微分すれば
$\nabla_{\dot\gamma}J(0)=w$ を得る．
Jacobi場の初期値一意性から指定の $J$ と一致する．
\end{proof}
\section{具体例}
\begin{example}[Euclid空間]
$\mathbb R^n$ の通常の計量では $R=0$ であり，
測地線に沿うJacobi場は $J(t)=a+tb$ の形となる．
\end{example}
\begin{proof}
Christoffel記号は零なので曲率も零である．
Jacobi方程式は各成分 $J''=0$ となり，
全解が一次式である．
\end{proof}

\section{到達課題の完全解答}
\noindent\textbf{定曲率空間のJacobi方程式．}
\begin{proposition}[定曲率の垂直Jacobi場]
単位速度測地線 $\gamma$ に沿い，
$\dot\gamma$ に垂直で平行な単位ベクトル場 $E$ を取る．
断面曲率が一定 $K$ なら，
$J(t)=j(t)E(t)$ がJacobi場であることと
\[
j''(t)+Kj(t)=0
\]
は同値である．
特に $j(0)=0$，$j'(0)=1$ の解は
\[
j(t)=
\begin{cases}
\sin(\sqrt Kt)/\sqrt K&K>0,\\
t&K=0,\\
\sinh(\sqrt{-K}t)/\sqrt{-K}&K<0
\end{cases}
\]
である．
\end{proposition}
\begin{proof}
$\nabla_{\dot\gamma}E=0$ より
$\nabla_{\dot\gamma}^2J=j''E$ となる．
定曲率の公式
$R(J,\dot\gamma)\dot\gamma=KJ$
をJacobi方程式に代入すると
$j''+Kj=0$ を得る．
表示した三つの関数はいずれも
対応する方程式と初期条件を満たす．
二階常微分方程式の一意性から結論が従う．
\end{proof}
\begin{corollary}[正曲率の場合の最初の零点]
$K>0$ の場合，
$j(t)$ は正の時刻 $\pi/\sqrt K$ で初めて零になる．
一方 $K\le0$ では $t>0$ に零点を持たない．
\end{corollary}
\begin{proof}
三つの明示式で直接確かめる．
\end{proof}
\noindent\textbf{到達確認．}
$K=4$ の場合，$j(0)=0$，$j'(0)=1$ を満たす
Jacobi場のスカラー係数と最初の正の零点を求めよ．
\par\noindent\textbf{解答．}
$j(t)=\sin(2t)/2$，最初の零点は $\pi/2$ である．
\end{document}
